\documentclass{article}
\usepackage[utf8]{inputenc}
\usepackage{geometry}
\geometry{a4paper, margin=1in}

\title{Análisis del Proyecto de Seguridad Informática}
\author{Gabriel}
\date{\today}

\begin{document}

\maketitle

\section{Resumen del Proyecto}
El proyecto es una herramienta de línea de comandos para auditoría de seguridad. Actualmente, 
implementa un módulo que realiza consultas DNS para los registros A y AAAA de un objetivo 
determinado.

\section{Estructura del Código}
\begin{itemize}
    \item \texttt{auditoria.py}: El script principal que maneja los argumentos de la línea de 
    comandos usando \texttt{argparse}.
    \item \texttt{modulos.py}: Contiene la lógica de la auditoría, en este caso, la función 
    \texttt{get\_dns\_info}.
    \item \texttt{requirements.txt}: Define las dependencias del proyecto, que son 
    \texttt{dnspython} y \texttt{argparse}.
\end{itemize}

\section{Puntos Fuertes}
\begin{itemize}
    \item \textbf{Buena Organización:} El código está bien estructurado, separando la lógica 
    principal de los módulos de funcionalidad, lo que facilita su mantenimiento y expansión.
    \item \textbf{Interfaz de CLI Clara:} El uso de \texttt{argparse} es adecuado para una 
    herramienta de línea de comandos.
    \item \textbf{Manejo de Errores Básico:} La función de DNS incluye un manejo de excepciones 
    para cuando no se encuentran registros.
\end{itemize}

\section{Áreas de Mejora}
\begin{itemize}
    \item \textbf{Ampliar Módulos Existentes:} La función de DNS podría extenderse para consultar 
    otros tipos de registros (ej. MX, NS, TXT, SOA).
    \item \textbf{Formato de Salida:} La información podría presentarse en un formato más 
    estructurado como JSON, lo que permitiría que otras herramientas procesen la salida 
    fácilmente.
    \item \textbf{Pruebas Unitarias:} Añadir pruebas para las funciones en \texttt{modulos.py} 
    aseguraría que el código funciona como se espera y previene regresiones.
\end{itemize}

\end{document}
